% [原文 p.191]
\chapter{极小模型纲领的奇点}

本章的目的是研究极小模型纲领中出现的高维奇点。

§5.1 研究 Cohen--Macaulay 奇点与有理奇点的性质。这些结果与其说属于极小模型纲领，不如说属于一般代数几何。本书只在自身目的所需的程度上展开这套理论。

§5.2 考察 dlt 对的局部性质。最重要的结果之一，是它们的有理性定理 5.22。由此，常常可以把 $\Q$-Cartier Weil 除子当作 Cartier 除子来对待（命题 5.26、推论 5.27）。

§5.3 用于详细研究三维终端奇点与标准奇点。首先，把它们与二维 Du Val 奇点、椭圆奇点联系起来（定理 5.34、定理 5.35）。随后，把第 4 章的结果应用于这些研究。三维终端奇点的完全分类在定理 5.43 中叙述，但本书不需要这一结果。

§5.4 汇集关于伴随与伴随反转的种种结果。借助伴随反转，常常可以把 dlt 对的问题归约到低维。主要结果是定理 5.50。类似的手法也用于研究 dlt 对边界的既约部分。这些结果在第 7 章中很重要。

§5.5 是 CM 层对偶理论的简短入门。

从本章起，将会遇到若干这样的定理：论断本身是局部的，证明却在射影簇
% [原文 p.192]
的情形下容易得多。最好的例子是定理 5.22。追溯困难的根源，多数情况下会归结到对偶理论：射影簇上的对偶理论，比固有态射的相对对偶理论容易叙述得多，证明也容易得多。

在这种情形下，定理按一般形式叙述，证明只处理射影簇的情形，对一般的代数情形或复解析情形则只给出引用文献。

\section{有理奇点}
\takeaway{本节复习 CM 层与有理奇点的基本性质。核心结果是定理 5.10 的三条等价刻画：存在有理消解、一切消解皆有理、$Y$ 为 CM 且 $f_*\omega_X=\omega_Y$。商奇点是最典型的例子。}

本节复习 CM 层与有理奇点的基本性质。除特别声明外，概形均假定为 Noether 的，层均取凝聚层。

本节的结果虽然是对代数簇叙述的，但全部在复解析情形下也同样成立。

\begin{definition}{5.1}
设 $(R,\mathfrak{m})$ 是 Noether 局部环，$N$ 是有限 $R$-模。记 $\dim N:=\dim\Supp N$。当下列等价条件之一成立时，称 $N$ 为 CM（Cohen--Macaulay 的缩写）。
\begin{enumerate}[label=(\arabic*),leftmargin=*]
\item 存在长度为 $r=\dim N$ 的 $N$-正则列（regular sequence）$x_1,\dots,x_r\in\mathfrak{m}$。即对每个 $i$，$x_i$ 是 $N/(x_1,\dots,x_{i-1})N$ 的非零因子。
\item 若 $x_1,\dots,x_r\in\mathfrak{m}$（$r=\dim N$）满足 $\dim N/(x_1,\dots,x_r)N=0$，则 $x_1,\dots,x_r$ 是 $N$-正则列。
\end{enumerate}
概形 $X$ 上的凝聚层 $F$，若对每个 $x\in X$，茎 $F_x$ 在 $\cO_{X,x}$ 上是 CM 的，则称 $F$ 为 CM。

概形 $X$ 的结构层 $\cO_X$ 为 CM 时，称 $X$ 为 CM。
\end{definition}

等价性 (1)$\Leftrightarrow$(2) 容易证明。例如参见 [Mat86, chap.6]。

\begin{definition}{5.2}
Noether 概形 $X$ 上的层 $F$，若对每个 $x\in X$，茎 $F_x$ 具有长度至少为 $\min\{d,\dim\cO_{X,x}\}$ 的正则列，则称 $F$ 为 $S_d$。因此，记 $d=\dim X$，则 $F$ 为 CM 当且仅当 $F$ 为 $S_d$。

一般地，凝聚层 $F$ 的对偶（dual）$F^*$ 由 $F^*:=\operatorname{Hom}_{\cO_X}(F,\cO_X)$ 定义，
% [原文 p.193]
于是有到双对偶的自然映射 $F\to(F^*)^*$。当它是同构映射时，称 $F$ 为反射的（reflexive）。

本书需要的只是 $S_1$ 与 $S_2$ 这两个条件。回顾：设 $F$ 为凝聚层，其相伴点（associated point）是指这样的点 $x\in X$——茎 $F_x$ 含有与 $x$ 的剩余域同构的子层。于是
\[
\begin{aligned}
F\ \text{为}\ S_1 &\iff F\ \text{的相伴点全为}\ X\ \text{的泛点}\\
&\iff F\ \text{无挠（torsion-free，当}\ X\ \text{不可约且既约时）},\\
F\ \text{为}\ S_2 &\iff F\ \text{为}\ S_1\ \text{且反射（当}\ X\ \text{正规时）}.
\end{aligned}
\]
\end{definition}

% [Tutor Note: 本书只用 $S_1$、$S_2$：$S_1$ 排除嵌入分量，$S_2$ 使无挠层与反射层重合。正规概形恰为 $S_2$ 且 $R_1$，这正是后文反复用 Hartogs 现象延拓迹与范数的根据。]

\begin{propositionp}{5.3}
设 $(R,\mathfrak{m})$ 是 Noether 局部环，$N$ 是有限 $R$-模。
\begin{enumerate}[label=(\arabic*),leftmargin=*]
\item 设 $x\in\mathfrak{m}$ 是 $N$ 上的非零因子。此时，$N$ 为 CM（或 $S_d$）当且仅当 $N/xN$ 为 CM（或 $S_{d-1}$）。
\item 设 $R$ 正则，$N$ 有限且 $\dim N=\dim R$。此时，$N$ 为 CM 当且仅当 $N$ 为自由 $R$-模。
\end{enumerate}
\end{propositionp}

\begin{proof}
这些证明都比较容易。(1) 见 [Mat86, 17.3]。(2) 是 [Mat86, 19.1] 的一个非常特殊的情形，但用关于 $\dim R$ 的归纳法直接证明更为容易。
\end{proof}

\begin{propositionp}{5.4}
设 $f\colon X\to Y$ 是 Noether 概形之间的有限态射，$F$ 是 $X$ 上的凝聚层，且 $\Supp F$ 为纯 $r$ 维。此时
\[
F\ \text{为 CM（或}\ S_d\text{）}\iff f_*F\ \text{为 CM（或}\ S_d\text{）}.
\]
\end{propositionp}

\begin{proof}
不妨假定 $(R,\mathfrak{m})$ 是 Noether 局部环，$Y=\Spec(R,\mathfrak{m})$，$Q$ 是 $Y$ 的闭点。记 $\{P_1,\dots,P_s\}=f^{-1}(Q)$。此时 $X$ 是由一个半局部环决定的仿射概形，该半局部环的极大理想恰为诸 $P_i$。选取 $y_1,\dots,y_r\in\mathfrak{m}$，使 $\dim F/(y_1,\dots,y_r)F=0$。此时
\[
\begin{aligned}
F\ \text{为 CM} &\iff \text{对一切}\ i,\ F_{P_i}\ \text{为 CM}\\
&\iff \text{对一切}\ i,\ y_1,\dots,y_r\ \text{是}\ F_{P_i}\text{-正则的}\\
&\iff y_1,\dots,y_r\ \text{是}\ F\text{-正则的}\\
&\iff y_1,\dots,y_r\ \text{是}\ f_*F\text{-正则的}.
\end{aligned}
\]
同样的论证对 $S_d$ 也有效。
\end{proof}

% [原文 p.194]
\begin{corollaryp}{5.5}
设 $f\colon X\to Y$ 是纯 $n$ 维概形之间的有限态射，$F$ 是 $X$ 上的凝聚层。此时，$F$ 在 $Y$ 上平坦当且仅当 $f_*F$ 局部自由。若进一步 $Y$ 正则，则 $f$ 平坦当且仅当 $X$ 为 CM。
\end{corollaryp}

\begin{proof}
由 [Mat69, 3.G]，$F$ 在 $Y$ 上平坦当且仅当 $f_*F$ 局部自由。若 $Y$ 正则，则这与“$F$ 为 CM”等价（命题 5.3(2)）。
\end{proof}

\begin{definition}{5.6}
设 $X$ 是概形，$A$ 是 $\cO_X$-代数的局部自由层。取开集 $U\subset X$，使 $A|_U\cong\cO_U^r$。由 $b\in H^0(U,A|_U)$ 作左乘，可以看作 $\operatorname{Hom}_U(\cO_U^r,\cO_U^r)$ 中的元素 $\varphi(b)$。

$\varphi(b)$ 的迹（或行列式）不依赖同构 $A|_U\cong\cO_U^r$ 的选取，故称为 $b$ 的迹（trace）（或范数（norm））。迹给出 $\cO_X$-同态 $\operatorname{Trace}\colon A\to\cO_X$，范数给出乘法映射 $\operatorname{Nm}\colon A\to\cO_X$。

若 $\cO_X$ 为 $S_2$，且 $\cO_X$-代数层 $B$ 在余维至少为 $2$ 的子集 $Z\subset X$ 之外局部自由，则 $\operatorname{Trace}_{X\setminus Z}\colon B|_{X\setminus Z}\to\cO_{X\setminus Z}$（或 $\operatorname{Nm}_{X\setminus Z}\colon B|_{X\setminus Z}\to\cO_{X\setminus Z}$）唯一地延拓为 $\operatorname{Trace}\colon B\to\cO_X$（或 $\operatorname{Nm}\colon B\to\cO_X$）。该延拓也称为迹（或范数）。

若有限态射 $g\colon Y\to X$ 在余维至少为 $2$ 的子集 $Z\subset X$ 之外平坦，且 $X$ 为 $S_2$，则取 $B=g_*\cO_Y$，便得到 $\operatorname{Trace}_{Y/X}\colon g_*\cO_Y\to\cO_X$ 与 $\operatorname{Nm}_{Y/X}\colon g_*\cO_Y\to\cO_X$。注意：若 $X$ 与 $Y$ 正规且 $g$ 满射，则这些假定自然满足。

为避免误解，补充一句：命题 5.77 中定义的迹，在某种意义上应称为此处所定义之迹的对偶。

也可以把迹看作双线性映射（内积）$B\times B\to\cO_X$，$(b,b')\mapsto\operatorname{Trace}(bb')$。若 $g$ 在开集 $U\subset X$ 上（即在原像 $g^{-1}(U)$ 上）平展，则该内积在 $U$ 上非退化。这可由引理 5.78 得到。
\end{definition}

% [Tutor Note: 迹与范数的延拓依赖 $X$ 为 $S_2$：余维至少 $2$ 的例外集挡不住截面延拓，故先在 $X\setminus Z$ 上定义，再唯一延拓过去。此处的正规性假定是实质条件，并非装饰。]

\begin{propositionp}{5.7}
设 $f\colon X\to Y$ 是纯 $n$ 维概形之间的有限满射。
\begin{enumerate}[label=(\arabic*),leftmargin=*]
\item 若 $X$ 为 CM，且 $\cO_Y$ 是 $f_*\cO_X$ 的直和因子，则 $Y$ 也为 CM。
\item 若 $Y$ 是域上的正规概形，$X$ 没有余维 $1$ 的相伴点，且特征不整除 $\deg f$，则 $\cO_Y$ 是 $f_*\cO_X$ 的直和因子。
\end{enumerate}
\end{propositionp}

\begin{proof}
若 $\cO_Y$ 是 $f_*\cO_X$ 的直和因子，则 $\cO_Y$ 中的序列为 $\cO_Y$-正则列与为 $f_*\cO_X$-正则列是等价的。由此得到 (1)。

若 $Y$ 正规，则 $\frac{1}{\deg f}\operatorname{Trace}_{X/Y}\colon f_*\cO_X\to\cO_Y$ 给出内射 $\cO_Y\to f_*\cO_X$ 的分裂。
\end{proof}

从高维双有理几何的立场看，CM 奇点这一类仍然过宽。有理奇点是更窄的一类。可以说，有理奇点就是不影响结构层与对偶化层的上同调性质的奇点类。

\begin{definition}{5.8}
设 $Y$ 是特征 $0$ 的域上的簇，$f\colon X\to Y$ 是奇点消解。当下列 $2$ 个条件成立时，称 $f\colon X\to Y$ 为有理消解（rational resolution）。
\begin{enumerate}[label=(\arabic*),leftmargin=*]
\item $f_*\cO_X=\cO_Y$（与 $Y$ 正规等价）。
\item $R^if_*\cO_X=0$（$\forall i>0$）。
\end{enumerate}
若一切奇点消解 $f\colon X\to Y$ 都是有理的，则称 $Y$ 具有有理奇点（rational singularity）。
\end{definition}

% [原文 p.195]
\par\medskip\noindent\textbf{5.9}\quad 有理奇点的概念在正特征下运作得并不好。为得到好的理论，还须假定 $R^if_*\omega_X=0$（$i>0$）。（特征 $0$ 时，该假定由推论 2.68 满足。）然而，正特征下的非奇异点是否满足这一假定，目前尚不清楚。\hfill$\square$

\begin{theorem*}{5.10}
设 $Y$ 是特征 $0$ 的域上的簇。则下列条件等价。
\begin{enumerate}[label=(\arabic*),leftmargin=*]
\item $Y$ 具有有理消解。
\item $Y$ 的所有奇点消解都是有理的。
\item $Y$ 为 CM，且对奇点消解 $f\colon X\to Y$ 有 $f_*\omega_X=\omega_Y$（即命题 5.77 中定义的映射 $\operatorname{Trace}_{X/Y}$ 为同构）。
\end{enumerate}
\end{theorem*}

% [Tutor Note: 定理 5.10 把“所有消解都有理”这一看似极强的定义化归为单个消解上的条件。关键一步是 $f_*\omega_X$ 不依赖消解的选取：任一消解都可被第三个消解支配。]

% [原文 p.196]
由定理 2.48，由此立即得到：

\begin{corollary}{5.11}
若 $X$ 是 $\C$ 上的簇，则
\[
X\ \text{具有有理奇点}\iff X^{an}\ \text{具有有理奇点}.
\]
\end{corollary}

\begin{proof}[定理 5.10 的证明]
设 $f\colon X\to Y$ 是奇点消解，$g\colon X'\to X$ 是 $X$ 的奇点消解。此时 $g_*\omega_{X'}=\omega_X$（命题 5.77(3)）。于是得到 $f_*\omega_X=(f\circ g)_*\omega_{X'}$。由于存在同时支配 $Y$ 的任意两个奇点消解的奇点消解，这说明 $f_*\omega_X$ 不依赖奇点消解 $f\colon X\to Y$ 的选取。因此条件 (3) 不依赖消解的选取，故由下述引理即证得定理。
\end{proof}

\begin{lemmap}{5.12}
设 $Y$ 是特征 $0$ 的域 $k$ 上 $n$ 维簇，$f\colon X\to Y$ 是奇点消解。则下列条件等价。
\begin{enumerate}[label=(\arabic*),leftmargin=*]
\item $f$ 是有理消解。
\item $Y$ 为 CM，且 $f_*\omega_X=\omega_Y$。
\end{enumerate}
\end{lemmap}

\begin{proof}
只处理 $Y$ 为射影簇的情形。一般情形下比较简单的证明见 [Kol97, Sec.11]。

设 $D$ 是 $Y$ 上的丰富 Cartier 除子。对 $i>0$、$r>0$ 有 $H^i(X,\omega_X(rf^*D))=0$（定理 2.64），故由 Serre 对偶得到
\begin{equation}
H^{n-i}(X,\cO_X(-rf^*D))=0\qquad(i>0,\ r>0).\tag{5.1}
\end{equation}
还要使用下述 Leray 谱序列：
\begin{equation}
E_2^{i,j}=H^i(Y,R^jf_*\cO_X(-rD))\Rightarrow H^{i+j}(X,\cO_X(-rf^*D)).\tag{5.2}
\end{equation}
先证明 (1)$\Rightarrow$(2)。由假定 $R^jf_*\cO_X=0$（$j>0$），用（5.2）可得 $H^j(Y,\cO_Y(-rD))\cong H^j(X,\cO_X(-rf^*D))$。由（5.1），这意味着 $H^j(Y,\cO_Y(-rD))=0$（$r>0$，$n>j$）。因此由推论 5.72，$Y$ 为 CM。$Y$ 上的 Serre 对偶对 CM 层成立（定理 5.71）。由 $j=n$ 时的上述同构得到
\[
h^0(Y,\omega_Y(rD))=h^0(X,\omega_X(rf^*D))=h^0(Y,f_*\omega_X(rD))\qquad(r>0).
\]
这意味着 $f_*\omega_X=\omega_Y$。

其次，用关于 $n$ 的归纳法证明 (2)$\Rightarrow$(1)。对所有 $i>0$，
% [原文 p.197]
证明 $R^if_*\cO_X=0$ 在某 $0$ 维集合之外成立。为此，取 $H$ 为 $Y$ 的一般超平面截面，记 $H':=f^{-1}H$。则 $f\colon H'\to H$ 是 CM 概形 $H$ 的奇点消解。注意对 $X$ 上的凝聚层 $F$ 有 $(R^if_*F)\otimes\cO_H=R^if_*(F\otimes\cO_H)$。这是因为可以选取 $H$，使它不含 $R^{i+1}f_*F$ 的任何相伴点。因此，在 $X$ 与 $Y$ 上使用命题 5.73，得到
\[
\begin{aligned}
f_*\omega_{H'} &= f_*\bigl(\omega_X(H')\otimes\cO_{H'}\bigr)\\
&= \cO_H(H)\otimes f_*\omega_X=\cO_H(H)\otimes\omega_Y=\omega_H.
\end{aligned}
\]
由归纳法，$\cO_H\otimes R^if_*\cO_X=R^if_*\cO_{H'}=0$ 对一切 $i>0$ 成立。这就证明了断言。

因此，当 $p,q>0$ 时 $H^p(Y,R^qf_*\cO_X(-rD))=0$。又当 $p<n$ 且 $q=0$ 时，只要 $r\gg0$，该等式也成立（用推论 5.72，由假定导出）。因此由谱序列（5.2）得到
\[
H^0(Y,R^qf_*\cO_X(-rD))=0\qquad(q<n-1),
\]
\[
H^0(Y,R^{n-1}f_*\cO_X(-rD))\cong\ker\bigl[H^n(\cO_Y(-rD))\to H^n(\cO_X(-rf^*D))\bigr].
\]
若 $q>0$，则 $R^qf_*\cO_X$ 的支集至多为 $0$ 维，故由第一式得到 $R^qf_*\cO_X=0$（$0<q<n-1$）。$\alpha$ 与下述映射对偶：
\[
H^0(Y,\omega_Y(rD))\to H^0(X,\omega_X(rf^*D))=H^0(Y,f_*\omega_X(rD)),
\]
而由 $f_*\omega_X=\omega_Y$，后者为同构。由此 $\alpha$ 为同构，$H^0(Y,R^{n-1}f_*\cO_X(-rD))=0$，从而得到 $R^{n-1}f_*\cO_X=0$。
\end{proof}

% [Tutor Note: $\alpha$ 即 p.197 显示中的映射 $H^n(\cO_Y(-rD))\to H^n(\cO_X(-rf^*D))$。它与 $H^0(Y,\omega_Y(rD))\to H^0(X,\omega_X(rf^*D))$ 对偶；后者由 $f_*\omega_X=\omega_Y$ 为同构，故 $\alpha$ 亦为同构。]

% [原文 p.198]
\begin{propositionp}{5.13}
设 $f\colon X\to Y$ 是特征 $0$ 的域上 $n$ 维簇之间的有限态射。假定 $X$ 只有有理奇点，$Y$ 正规。此时 $Y$ 也只有有理奇点。
\end{propositionp}

\begin{proof}
由命题 5.7 已知 $Y$ 为 CM。设 $g^Y\colon Y'\to Y$ 是奇点消解。由此得到下述交换图。
\[
\begin{tikzcd}
X' \arrow[r,"g^X"] \arrow[d,"f'"] & X \arrow[d,"f"] \\
Y' \arrow[r,"g^Y"] & Y
\end{tikzcd}
\]
其中 $f'$ 有限，$X'$ 正规。考虑 $X'$ 的奇点消解，由 $X$ 只有有理奇点即可得到 $g^X_*\omega_{X'}=\omega_X$。再考察下述交换图。
\[
\begin{tikzcd}
f_*g^X_*\omega_{X'} \arrow[r,phantom,"="] \arrow[d,equal] & f_*\omega_X \arrow[r,"\operatorname{Trace}_{X/Y}"] & \omega_Y \arrow[d,equal] \\
g^Y_*f'_*\omega_{X'} \arrow[r,"\operatorname{Trace}_{X'/Y'}"] & g^Y_*\omega_{Y'} \arrow[r,"c"] & \omega_Y
\end{tikzcd}
\]
首先，$\cO_Y$ 是 $f_*\cO_X$ 的直和因子。因此由命题 5.68，$\omega_Y$ 是 $f_*\omega_X$ 的直和因子，而命题 5.77 中定义的 $\operatorname{Trace}_{X/Y}$ 是满射。由此 $c$ 为满射，故 $Y$ 只有有理奇点。
\end{proof}

商奇点是值得关注的有理奇点之例。

\begin{definition}{5.14}
设 $x\in X$ 是复解析空间的芽。若存在非奇异的芽 $0\in Y$ 与作用在 $0\in Y$ 上的有限群 $G$，使得 $(x\in X)\cong(0\in Y)/G$，则称 $X$ 具有商奇点（quotient singularity）。

设 $X$ 是 $\C$ 上的代数簇。若 $X^{an}$ 在各点的芽都只有商奇点，则称 $X$ 只有商奇点（或简称具有商奇点）。
\end{definition}

\begin{propositionp}{5.15}
设 $X$ 是 $\C$ 上的代数簇或解析空间，且只有商奇点。此时 $X$ 具有有理奇点，且 $X$ 是 $\Q$-分解的。
\end{propositionp}

\begin{proof}
由推论 5.11 与命题 5.13，商奇点是有理的。

设 $x\in X$ 是闭点。$\cO_{x,X}$ 是其局部环，$\cO^{an}_{x,X}$ 是其解析局部环，$\hat{\cO}_{x,X}$ 是二者的共同完备化。由 [Mat69, 24.E]，对理想 $I\subset\cO_{x,X}$，$I$ 为主理想当且仅当 $I\cO^{an}_{x,X}$ 为主理想，当且仅当 $I\hat{\cO}_{x,X}$ 为主理想。因此只需证明解析商奇点是 $\Q$-分解的。这由下述引理得到。
\end{proof}

% [Tutor Note: 商奇点既是有理奇点，又是 $\Q$-分解奇点，属于最温和的奇点之列。证明分两步：有理性由有限商保持有理性得到（命题 5.13），$\Q$-分解性由引理 5.16 得到。]

% [原文 p.199]
\begin{lemmap}{5.16}
设 $f\colon X\to Y$ 是正规簇之间的有限满射。若 $X$ 是 $\Q$-分解的，则 $Y$ 也是 $\Q$-分解的。
\end{lemmap}

\begin{proof}
设 $F$ 是 $X$ 上任一素除子。要证 $f(F)$ 在 $X$ 上为 $\Q$-Cartier。存在 $a\in\Z_{>0}$，使 $aF$ 为 Cartier。对任意的 $x\in X$，存在开集 $U\ni x$，使 $aF$ 在 $V:=f^{-1}(U)$ 上由一个方程 $\phi=0$ 定义。此时，按范数的构造，$\operatorname{Nm}_{V/U}(\phi)=0$ 在 $U$ 上定义了子集 $f(F)$。因此，$X$ 上的素除子 $f(F)$ 是 $\Q$-Cartier 的。
\end{proof}

% [Tutor Note: 原文两处把 $f(F)$ 所在的空间写作 $X$，按前后文应为 $Y$：$f(F)$ 是 $Y$ 上的素除子，$F$ 才是 $X$ 上的。底本如此，照译并注明。]

\section{对数终端奇点}
\takeaway{本节研究 dlt 对的局部性质。核心结论是 dlt 对具有有理奇点（定理 5.22），从而常可把 $\Q$-Cartier Weil 除子当作 Cartier 除子使用。基本工具是超平面截面与指数 $1$ 循环覆盖。}

本节的目标是详细研究对数终端奇点。在一般维数下成立的结果，将注明其适用范围。

与上一节一样，全部按代数簇的设定来叙述，但在解析空间的情形下也同样成立。

从两种基本的归约方法开始：一是超平面截面，二是取循环覆盖。

\begin{lemmap}{5.17}
设 $(X,\Delta)$ 是对，$|H|$ 是 $X$ 上的无基点线性系，$H_g\in|H|$ 是一般元。则下列成立。
\begin{enumerate}[label=(\arabic*),leftmargin=*]
\item $\operatorname{discrep}(X,\Delta)\le\operatorname{discrep}(H_g,\Delta|_{H_g})$。
\item $\operatorname{discrep}(X,\Delta+H_g)=\min\{0,\operatorname{discrep}(X,\Delta)\}$。
\end{enumerate}
\end{lemmap}

\begin{proof}
设 $f\colon X'\to X$ 是 $(X,\Delta)$ 的对数奇点消解，记 $H'_g:=f^{-1}(H_g)$。则 $f$ 也是 $(X,\Delta+H_g)$ 的对数奇点消解，且 $f\colon H'_g\to H_g$ 是 $(H_g,\Delta|_{H_g})$ 的对数奇点消解。(1) 由伴随公式得到，(2) 由推论 2.32 成立。
\end{proof}

下述结果用于描述终端奇点与标准奇点在余维 $2$ 处的行为。

\begin{corollaryp}{5.18}
若 $X$ 终端，则 $X$ 在余维 $2$ 处非奇异。若 $X$ 标准，则 $K_X$ 在余维 $2$ 处为 Cartier。
\end{corollaryp}

% [原文 p.200]
这里，若条件 $P$ 在余维至多为 $k$ 的任意点处成立，则称 $P$ 在余维 $k$ 处成立（in codimension $k$）。特别地，当 $P$ 为开条件时，这等价于说 $P$ 不成立的点全体是余维至少为 $k+1$ 的闭集。

\begin{proof}
把引理 5.17 应用 $(\dim X-2)$ 次，归结为二维终端奇点、标准奇点的描述（定理 4.5）即可。
\end{proof}

% [Tutor Note: 这是本章反复使用的降维范式：对一般超平面截面反复用引理 5.17，把高维断言化为曲面上第 4 章的已知结论。终端性与标准性正是这样逐维继承下来的。]

其次，说明指数 $1$ 覆盖的手法。它频繁用于把问题归约到 $K_X$ 为 Cartier 的情形。

\begin{definition}{5.19}
设 $X$ 是正规簇，$D$ 是 $X$ 上的 $\Q$-Cartier Weil 除子。使 $rD$ 为 Cartier 的最小正整数 $r$，称为 $D$ 在 $X$ 上的指数（index）。特别地，$D$ 的指数为 $1$ 当且仅当 $D$ 为 Cartier。

$K_X$ 的指数也称为 $X$ 的指数（index）。

假定 $\cO_X(rD)\cong\cO_X$（对 $X$ 的任一点，把 $X$ 限制到其适当的邻域上，此式即成立），并选取处处非零的截面 $s\in H^0(X,\cO_X(rD))$。设 $p\colon Z\to X$ 是对应的循环覆盖（定义 2.50）。$p$ 在 $X\setminus\Sing X$ 上平展。特别地，$K_Z=p^*K_X$ 成立。$p^*D$ 是 Cartier 除子，且由引理 2.53，$p^*D\sim0$ 成立。

若 $D=K_X$，则称 $p\colon Z\to X$ 为 $X$ 的指数 $1$ 覆盖（index 1 cover）。

$s$ 的选取并不重要，故也常把 $Z$ 当作唯一确定的指数 $1$ 覆盖来处理。
\end{definition}

% [Tutor Note: 指数 $1$ 覆盖把 $K_X$ 非 Cartier 的问题搬到 $K_Z$ 为 Cartier 的 $Z$ 上。由于 $p$ 在奇点外平展，$K_Z=p^*K_X$；再用命题 5.20 比较两边的差异系数即可。]

循环覆盖之所以有用，是因为存在比较有限态射下差异系数的一般原理。这一类结果最早出现在 [Rei80] 中。

\begin{propositionp}{5.20}
设 $g\colon X'\to X$ 是 $n$ 维正规簇之间的有限态射。设 $\Delta=\sum a_iD_i$ 是 $X$ 上的 $\Q$-除子，$\Delta':=\sum a'_jD'_j$ 是 $X'$ 上的 $\Q$-除子，且满足 $K_{X'}+\Delta'=g^*(K_X+\Delta)$。则下列成立。
\begin{enumerate}[label=(\arabic*),leftmargin=*]
\item $K_X+\Delta$ 为 $\Q$-Cartier 当且仅当 $K_{X'}+\Delta'$ 为 $\Q$-Cartier。
\item $(\deg g)\bigl(\operatorname{discrep}(X,\Delta)+1\bigr)\ge(\operatorname{discrep}(X',\Delta')+1)$。
\item $\operatorname{discrep}(X',\Delta')\ge\operatorname{discrep}(X,\Delta)$。
\item $(X,\Delta)$ 为 klt（或 lc）当且仅当 $(X',\Delta')$ 为 klt（或 lc）。
\end{enumerate}
\end{propositionp}
